\documentclass[pdflatex,sn-chicago]{sn-jnl}

\usepackage{graphicx}
\usepackage{booktabs}
\usepackage{amsmath}
\usepackage{amssymb}
\usepackage{array}
\usepackage{caption}
\usepackage{hyperref}
\usepackage{url}

\title[Supplementary Information]{Supplementary Information for:\\A Training-Pool Replay Audit of the Released SLRTP2025 Back-Translation Evaluator: Non-Replication Across Independently Trained Checkpoints and a Readout-With-Generalization Signature}

\author*[1]{} % Authors retained for blind review
\affil*[1]{}

\begin{document}
\maketitle

\section*{Overview}
This supplement provides detailed appendices referenced in the main manuscript. Section labels match the labels used in the main text (e.g., \texttt{sec:appendix\_missing\_id} $\to$ Sup.\ Section~\ref{sec:appendix_missing_id}).


\section{Missing Test ID Sensitivity}
\label{sec:appendix_missing_id}
The official PHOENIX-2014T metadata contain 7,096/519/642 train/development/test IDs. The released SLRTP2025 pose records available to this study contain 7,060/515/641 sequences: 36/4/1 official IDs are absent from the released pose materialization. Their IDs and the reason \texttt{absent\_from\_slrtp2025\_released\_pose\_records} are enumerated in \texttt{artifact/manifests/exclusions.jsonl}. The missing test ID is \texttt{13December\_2010\_Monday\_tagesschau-6940}; its text/gloss metadata are not present in the released pose record, so length and template density cannot be reconstructed without an additional official annotation source. We therefore condition every result on the complete 641-record release rather than assert that the missing item is ignorable. As an empirical observed-extreme sensitivity scenario, assigning the missing item the per-item score of the largest PURE-over-GT outlier observed among the 641 released items changes the corpus-level gap by less than $\pm 0.2$ BLEU points, well within the reported 95\% paired-bootstrap interval $[+9.6,+12.4]$. This is an observed-extreme scenario, not a worst-case bound: corpus BLEU is non-linear and the missing item's reference, hypothesis, and n-gram sufficient statistics are unknown, so the true mathematical worst case is not covered. All primary conclusions are therefore restricted to the released 641-record materialization.

\section{Per-seed Extension Cells (707--1405)}
\label{sec:appendix_extension_cells}
The eight extension seeds are post-hoc additions that enter only the seed-level prediction interval and descriptive ranges, never the six primary comparisons. Per-seed PHX-public values (sacreBLEU 0--100 scale, canonical \S\ref{sec:retrieval} donor registry):

\begin{table}[h]
\centering
\scriptsize
\caption{Per-seed extension cells for the eight post-hoc extension seeds (707--1405). All values are PHX-public sacreBLEU on the 0--100 scale with the canonical \S\ref{sec:retrieval} donor registry (exact-normalized-text exclusion applied). All gaps are negative; these seeds enter only the prediction interval.}
\label{tab:extension_cells}
\begin{tabular}{lrrrr}
\toprule
Seed & dev BLEU-4 & GT & PURE & Gap \\
\midrule
707  & 6.69 & 9.86 & 8.97 & $-0.89$ \\
808  & 7.30 & 9.87 & 8.38 & $-1.48$ \\
909  & 8.96 & 9.20 & 8.31 & $-0.89$ \\
1001 & 9.07 & 8.81 & 7.95 & $-0.87$ \\
1102 & 7.61 & 9.75 & 8.63 & $-1.12$ \\
1203 & 8.14 & 9.06 & 8.11 & $-0.95$ \\
1304 & 8.18 & 9.68 & 8.50 & $-1.18$ \\
1405 & 8.12 & 8.05 & 6.86 & $-1.19$ \\
\bottomrule
\end{tabular}
\end{table}

\section{Analysis Provenance and Multiplicity}
\label{sec:provenance}
Table~\ref{tab:provenance} classifies every analysis in this manuscript by when and why it entered the paper, its evidentiary status, and its multiplicity handling. The single primary estimand is the PURE--GT decoded-sacreBLEU gap under the released evaluator (\S\ref{sec:method}); everything else is secondary, and secondary results are either Holm-corrected within a defined family or explicitly descriptive.

\begin{table}[ht]
\centering
\scriptsize\setlength{\tabcolsep}{2pt}
\caption{Analysis provenance. ``Prespecified'' = fixed before the headline result was observed; ``Post-hoc'' = added during the audit; ``Correction'' = supersedes an earlier pipeline error. C = confirmatory (frozen protocol, error control), D = descriptive (no formal error control).}
\label{tab:provenance}
\begin{tabular}{p{6.2cm}p{2.8cm}p{0.9cm}p{3.0cm}}
\toprule
Analysis & Provenance & Status & Multiplicity handling \\
\midrule
PURE--GT gap, original evaluator (primary estimand) & Exploratory-observed; protocol frozen afterwards & D & descriptive MC estimate $p\leq10^{-4}$ \\
Protocol checks (scorer, FPS, decoding sweep) & Prespecified & C/D & exact-match / descriptive \\
Conditioning controls (retrievers, gloss, random) & Prespecified & D & descriptive \\
14 reconstruction gaps + prediction interval & 6 prespecified; 8 post-hoc & C/D & DiD CIs; PI descriptive \\
Cross-metric gap table & Correction (metric alignment) & D & direction/sign only \\
Template-density strata; donor-exclusion $\tau$-sweeps & Prespecified & D & descriptive \\
Template-holdout, BT-retrain, cross-fit & Prespecified & C & paired bootstrap CIs \\
Donor-pool substitution; canonical B/C/D path & Post-hoc & D & query-index bootstrap CIs \\
Alternative paths, Shapley, S4 system; balance/caliper/signer & Post-hoc & D & path-sensitivity display \\
605-query factorial (cells, effects, interaction) & Prespecified design & C & Holm across 3 contrasts \\
Reference replacement (single/dual); overlap strata & Post-hoc & D & paired bootstrap of gap difference \\
Transcript-slot diagnostic + extractor audit & Post-hoc & D & descriptive; extractor P/R \\
Five related diagnostic views (\S\ref{sec:positive_evidence}) & Post-hoc & D & triangulation, not independent \\
Oracle/white-box transfer matrix & Post-hoc; appendix-only & D & permutation check \\
Competence gate, trajectories, Pareto, rescue (20); gate threshold sensitivity & Gate prespecified; rescue expanded & D & gate counts; 12-threshold grid \\
Big-arch variants; long-schedule runs & Post-hoc & D & competence attempts \\
In-sample readout diagnosis (train/dev/test) & Author-initiated & D & descriptive; no identity claim beyond released train.pt \\
Rank stability (9 systems $\times$ 7 evaluators) & Post-hoc & D & $\tau_b$/$\rho$/flip means \\
IPW pool-origin robustness & Post-hoc & D & balance robustness check \\
Slot hypothesis-side audit & Post-hoc & D & extractor P/R by system \\
Frozen-protocol confirmation (seeds 1506/1607) & Post-hoc & C (frozen protocol) & confirmation of family-side pattern \\
Knowledge-distillation competence ladder (15 planned; 15 evaluated) & Post-hoc & D & 15 with gap + dev \\
Per-$\alpha$ distillation breakdown (Q2) & Post-hoc & D & descriptive; all $\alpha$ firmly negative \\
Competence extrapolation (OLS/GP) & Post-hoc & D & descriptive; non-IID; LOO CV (5 folds) \\
Asymmetric floor-effect control (M4) & Post-hoc & D & BLEU-4/BLEU-1/chrF drops \\
Evaluator$\times$reference factorial (7$\times$3) & Post-hoc & D & descriptive decomposition \\
CSL-Daily readout-ratio check & Post-hoc & D & second-dataset boundary \\
Per-seed extension cells (8 seeds) & Post-hoc & D & descriptive; PI only \\
Ethics/licensing declarations & Updated each revision & --- & --- \\
\bottomrule
\end{tabular}
\end{table}

\section{Checkpoint Registry}
\label{sec:checkpoint_registry}
Table~\ref{tab:checkpoint_registry} is the definitive registry of every BT evaluator checkpoint used in this manuscript, with the analyses each family enters; all checkpoint counts in the text refer to these families.

\begin{table}[h]
\centering
\scriptsize\setlength{\tabcolsep}{3pt}
\caption{Canonical checkpoint registry. 70 planned training runs, all 70 completed training and beam-3 PHX-public evaluation. 67 non-holdout runs have dev metrics; 52 of these are non-distillation and enter the competence gate. 43 runs have both dev BLEU-4 and PHX-public PURE--GT gap (extrapolation-eligible). Plus the released original. Each family row's columns sum to the Total row.}
\label{tab:checkpoint_registry}
\begin{tabular}{p{4.2cm}p{1.0cm}p{1.5cm}p{1.4cm}p{1.4cm}p{3.0cm}}
\toprule
Family & Planned & Trained & Has gap (beam-3) & Has dev & Enters \\
\midrule
Reconstructions (primary) & 6 & 6 & 6 & 6 & multiseed, bootstrap, dose--resp. \\
Reconstructions (extension) & 8 & 8 & 8 & 8 & prediction interval, dose--resp. \\
Rescue lr (A/B $\times$ 4) & 8 & 8 & 0 & 8 & gate counts only \\
Rescue expanded & 12 & 12 & 1 & 12 & gate; wd0 also dose--resp., perturb. \\
Train-pool ladder & 4 & 4 & 4 & 4 & dose--resp., gap ranges \\
Config-faithful & 4 & 4 & 4 & 4 & dose--resp. \\
Step-faithful & 2 & 2 & 2 & 2 & dose--resp., validation-frequency \\
Large-arch (A1--A4) & 4 & 4 & 0 & 4 & competence attempt \\
Confirmation & 2 & 2 & 2 & 2 & dose--resp., family-side confirm. \\
Long-schedule & 2 & 2 & 1 & 2 & gate; best seed dose--resp. (dev 10.02) \\
BT-retrained & 1 & 1 & 0 & 0 & holdout series \\
Cross-fit A/B & 2 & 2 & 0 & 0 & holdout series \\
Distillation students & 15 & 15 & 15 & 15 & ladder, dose--resp. \\
\midrule
Total & 70 & 70 & 43 & 67 & \\
Released original & --- & --- & 1 & 1 & audit target (not trained) \\
\bottomrule
\end{tabular}
\par\smallskip
{\footnotesize Per-family columns sum to the Total row (e.g., Has gap: $6{+}8{+}0{+}1{+}4{+}4{+}2{+}0{+}2{+}1{+}0{+}0{+}15 = 43$). Gate-eligible = 52 non-distillation runs with dev metrics. Extrapolation-eligible = 43 runs with both dev BLEU-4 and PHX-public PURE--GT gap; the dose--response figure caption uses the same decomposition (14 reconstructions + 4 ladder + 4 config-faithful + 2 step-faithful + 2 confirmation + 1 long-schedule + 15 distillation + 1 wd0 rescue). All 15 distillation students completed training and beam-3 evaluation. The 1 wd0 rescue with PHX-public gap is the best rescue checkpoint (weight-decay 0, seed 202); the other 11 rescue-expanded runs have dev metrics but were not decoded on PHX-public. The second long-schedule seed was trained but not decoded on PHX-public.}
\end{table}

\section{Donor-Pool Substitution Decomposition Details}
\label{sec:appendix_decomposition}
This appendix gives the full path-sensitivity decomposition, factorial design, and balance checks referenced in \S\ref{sec:exposure_results}.

\paragraph{Path sensitivity.}
With SEEN-RAND640-MATCHED-v1 added, the systems form a $2\times2$ within-train design (pool size $\{$7060, 640$\}$ $\times$ selection $\{$max-Jaccard, Jaccard-matched$\}$) plus the test-pool cell. Because corpus BLEU is non-linear in the support, all contrasts use the common 622-query support. Every ordering telescopes to the same total $+14.7$ $[+13.3,+16.2]$, but the attribution varies by up to 7 BLEU points across orderings (Table~\ref{tab:path_sensitivity}). We report these as path-dependent descriptive contrasts, not mechanism decompositions.

\begin{table}[h]
\centering
\scriptsize
\caption{Path sensitivity of the donor-pool substitution decomposition under the original evaluator (common 622-query support; 10{,}000-resample query-index bootstrap CIs). S1=SEEN-PURE (7{,}060-pool, max-Jaccard), S2=SEEN-RAND640 (640-pool, max), S3=SEEN-MATCHED (7{,}060-pool, Jaccard-matched), S4=SEEN-RAND640-MATCHED (640-pool, matched), U=UNSEEN-PURE (test pool). Every ordering telescopes to the same total (A)~$=$ S1$-$U~$=+14.7$.}
\label{tab:path_sensitivity}
\begin{tabular}{llrrr}
\toprule
Ordering & Step & Size contrast & Selection contrast & Origin contrast \\
\midrule
Canonical (prespecified) & S1$\to$S2$\to$S3$\to$U & $+7.2$ $[+5.9,+8.4]$ & $-0.8$ $[-1.8,+0.2]$ & $+8.4$ $[+7.3,+9.4]$ \\
Alternative 1 & S1$\to$S2$\to$S4$\to$U & $+7.2$ $[+5.9,+8.4]$ & $+1.8$ $[+1.1,+2.5]$ & $+5.8$ $[+4.7,+6.8]$ \\
Alternative 2 & S1$\to$S3$\to$S4$\to$U & $+2.6$ $[+1.7,+3.5]$ & $+6.4$ $[+5.1,+7.7]$ & $+5.8$ $[+4.7,+6.8]$ \\
\midrule
Shapley average & over orderings & $+4.9$ $[+4.0,+5.8]$ & $+4.1$ $[+3.3,+4.9]$ & $+5.8$ $[+4.7,+6.8]$ \\
\bottomrule
\end{tabular}
\end{table}

\paragraph{Bootstrap family.}
All decomposition CIs use 10{,}000 resamples. Four resampling designs (query-index paired bootstrap; three one-way cluster bootstraps for the origin contrast on template-family, SEEN-donor, and UNSEEN-donor clusters; and a joint multiway replicate) give origin-contrast CIs of $[+7.2,+9.5]$ / $[+7.1,+9.7]$ / $[+7.1,+9.7]$ / $[+7.7,+9.1]$, all matching the query-index CI $[+7.3,+9.4]$ to within $\pm0.3$, so dependence through shared donors is negligible.

\paragraph{Factorial design.}
The 605-query factorial retains 605 of 641 targets (36 excluded for no joint match within tolerances; selection frozen before scoring). Table~\ref{tab:factorial_design} reports the balance. Table~\ref{tab:factorial_full} reports the full $2\times2$ cells, main effects, and interaction.

\begin{table}[h]
\centering
\small
\caption{Factorial design and balance (605 targets). SMD = standardized mean difference, seen vs unseen within stratum.}
\label{tab:factorial_design}
\begin{tabular}{lrrrr}
\toprule
Quantity & seen\_high & seen\_low & unseen\_high & unseen\_low \\
\midrule
LaBSE similarity (mean) & 0.813 & 0.654 & 0.776 & 0.622 \\
Duration (frames, mean) & 100.2 & 99.5 & 99.2 & 98.8 \\
Word count (mean) & 13.3 & 13.6 & 13.3 & 13.4 \\
Same-signer rate & 100\% & 100\% & 100\% & 100\% \\
\midrule
Similarity SMD (seen $-$ unseen) & \multicolumn{2}{r}{$+0.41$ (high)} & \multicolumn{2}{r}{$+0.38$ (low)} \\
\bottomrule
\end{tabular}
\end{table}

\begin{table}[h]
\centering
\scriptsize\setlength{\tabcolsep}{3pt}
\caption{Full $2\times2$ exposure factorial (605 targets; corpus sacreBLEU in BLEU points; NLL per token). CIs: 10{,}000-resample bootstrap (original) and seed-level $t$ across six reconstructions; Holm correction across three contrasts.}
\label{tab:factorial_full}
\begin{tabular}{llrrrrrrr}
\toprule
Metric & Evaluator & seen\_hi & seen\_lo & unseen\_hi & unseen\_lo & seen main & sim.\ main & inter. \\
\midrule
BLEU & Original & 11.12 & 2.31 & 7.19 & 2.08 & $\mathbf{+2.08}$ & $+6.96$ & $\mathbf{+3.70}$ \\
 &  &  &  &  &  & $[+1.6,+2.6]$ & $[+5.9,+8.1]$ & $[+2.6,+4.8]$ \\
BLEU & Seeds (mean) &  &  &  &  & $+0.51$ & $+4.06$ & $-0.39$ \\
\midrule
NLL & Original & 3.74 & 4.30 & 3.93 & 4.37 & $-0.128$ & $-0.498$ & $-0.124$ \\
NLL & Seeds (mean) &  &  &  &  & $-0.064$ & $-0.179$ & $-0.004$ \\
\bottomrule
\end{tabular}
\end{table}

\paragraph{Multivariate balance and robustness.}
On the 622-query common support, all six donor covariate SMDs between SEEN-PURE-MATCHED and UNSEEN-PURE are $|SMD|<0.12$ (Table~\ref{tab:multivariate_balance}). Caliper sensitivity is small (tolerance 0.05/0.10/0.20 gives $+8.2$/$+8.4$/$+8.6$). Signer-strict matching (tolerance 0.10, $n{=}580$) attenuates the residual to $+7.0$ $[+5.9,+8.1]$ --- sign and significance unchanged. Inverse propensity weighting drives all four covariate SMDs to 0.00 and leaves the residual at $+8.5$ $[+7.1,+10.1]$. The path-based estimate is robust to observed-covariate balancing but remains an association, not an identified exposure effect.

\begin{table}[h]
\centering
\scriptsize\setlength{\tabcolsep}{3pt}
\caption{Multivariate donor balance (622-query common support). All $|SMD|<0.12$.}
\label{tab:multivariate_balance}
\begin{tabular}{lrrrr}
\toprule
Covariate & SEEN-MATCHED mean & UNSEEN mean & SMD & Per-query overlap \\
\midrule
Source-text Jaccard & 0.411 & 0.415 & $-0.021$ & $|{\Delta}J|{\leq}0.05$: 96\% \\
LaBSE cosine & 0.782 & 0.774 & $+0.068$ & --- \\
Duration (s) & 3.76 & 3.74 & $+0.014$ & $|{\Delta}|\leq1$\,s: 57\% \\
Word count & 12.29 & 12.31 & $-0.003$ & $|{\Delta}|\leq2$: 69\% \\
Template density & 0.310 & 0.326 & $-0.075$ & --- \\
Same signer as query & 0.183 & 0.228 & $-0.111$ & both same-signer: 4.7\% \\
\bottomrule
\end{tabular}
\end{table}

\section{Config-Faithful and Perturbation Details}
\label{sec:appendix_config_faithful}

\paragraph{Epoch-validation approximation.}
The released \texttt{config.yaml} specifies \texttt{validation\_freq=14}, which JoeyNMT 1.2 interprets as 14 optimizer \emph{steps} ($\approx$28 steps/epoch means validating twice per epoch). Our initial implementation validated every 14 \emph{epochs} (an epoch-validation approximation); four seeds under this approximation all early-stop below the competence gate (dev BLEU-4 8.6--9.5), show negative PURE--GT gaps ($-1.1$ to $-0.2$), and family-range pass-through (1.21--1.24) and holdout (7.4--7.6) --- none reproduces the released signature (Table~\ref{tab:config_faithful}). Scaling capacity up (four larger variants: 4--6 layers, 384--512 hidden) overfits faster and reaches only dev BLEU 6.8--8.0. A long-schedule continuation (800 epochs, two seeds) pushes competence to a family maximum of 10.02 yet still shows no reversal (gap $-0.45$). The step-faithful re-run (\S\ref{sec:stepfaithful}) confirms the reconstruction plateaus below the released trajectory at matched step counts.

\begin{table}[h]
\centering
\scriptsize\setlength{\tabcolsep}{3pt}
\caption{Epoch-validation approximation (validate every 14 epochs, not steps): 300 epochs, dev-NLL selection. All four seeds produce negative PURE--GT gaps under the canonical \S\ref{sec:retrieval} donor registry; none shows the reversal or readout signature.}
\label{tab:config_faithful}
\begin{tabular}{lrrrr}
\toprule
Seed & ep. & dev NLL & GT & PURE gap \\
\midrule
101 & 27 & 2.11 & 10.12 & $-0.60$ \\
202 & --- & --- & 10.14 & $-1.70$ \\
303 & --- & --- & 8.74 & $-1.51$ \\
404 & --- & --- & 8.43 & $-0.52$ \\
\midrule
Released evaluator & 1{,}820 & --- & 12.78 & $+10.24$ \\
\bottomrule
\end{tabular}
\end{table}

\paragraph{Perturbation battery.}
\label{sec:appendix_perturbation}
Table~\ref{tab:perturbation} reports the full pose-perturbation battery. Reversing all frames collapses PURE from 23.79 to 7.72 (differential $16.1$ $[14.7,17.5]$; GT $5.5$ $[4.6,6.5]$); window-shuffling within one-second windows costs PURE $5.4$ $[4.4,6.6]$; 5\% spatial jitter is free ($-0.1$); masking hand or face joints collapses both systems. The best rescue checkpoint (wd0, seed 202) loses only $-4.3$ under reversal, confirming it is far less sensitive to donor temporal order.

\begin{table}[h]
\centering
\scriptsize\setlength{\tabcolsep}{3pt}
\caption{Pose-perturbation battery (corpus sacreBLEU, BLEU points) under the released evaluator and the best rescue checkpoint (wd0, seed 202).}
\label{tab:perturbation}
\begin{tabular}{llrrrrrr}
\toprule
Evaluator & System & identity & reversal & window shuffle & jitter 5\% & hand mask & face mask \\
\midrule
Original & GT-v1 & 12.78 & 7.25 & 11.97 & 12.51 & 1.35 & 3.15 \\
Original & TN-PURE-v1 & 23.79 & 7.72 & 18.35 & 23.91 & 1.30 & 4.06 \\
wd0 (seed 202) & GT-v1 & 10.73 & 6.26 & 10.54 & 10.60 & 0.63 & 1.58 \\
wd0 (seed 202) & TN-PURE-v1 & 10.98 & 6.66 & 9.65 & 10.92 & 0.58 & 1.43 \\
\bottomrule
\end{tabular}
\end{table}

\section{Cluster-Aware Bootstrap for Headline and Reference Analyses}
\label{sec:appendix_cluster_bootstrap}
The headline gap ($+10.24$, canonical \S\ref{sec:retrieval} registry) and the reference-attenuation result are reported with query-level paired bootstrap CIs. Because TN-PURE reuses donors and PHOENIX-2014T has template families and multiple signers, query-level resampling may underestimate uncertainty if dependence through shared donors, signers, or broadcast structure is non-negligible. Table~\ref{tab:cluster_bootstrap} reports interpretable cluster-aware resampling designs alongside the query-level baseline, using precomputed per-item n-gram sufficient statistics for 10{,}000 resamples each.

\begin{table}[h]
\centering
\scriptsize
\caption{Cluster-aware bootstrap for the headline PURE--GT gap (canonical \S\ref{sec:retrieval} registry, original refs). All CIs are 10{,}000-resample sacreBLEU \texttt{tok:13a} exp-smooth. With so few clusters (especially show with only 2), these are \emph{exploratory sensitivity checks}, not confirmatory tests.}
\label{tab:cluster_bootstrap}
\begin{tabular}{lrrrr}
\toprule
Analysis & Design & Clusters & Gap & 95\% CI \\
\midrule
Orig-ref gap (canonical) & Query & 641 & $+10.24$ & $[+8.89, +11.58]$ \\
& Signer & 9 & --- & --- \\
& Show & 2 & --- & --- \\
& Date & 297 & --- & --- \\
& Signer$\times$Show & 17 & --- & --- \\
\bottomrule
\end{tabular}
\end{table}

The signer-cluster CI ($[+9.08, +10.84]$, 9 clusters) is slightly tighter than the query-level CI ($[+8.45, +11.40]$), consistent with the expectation that within-signer items are positively correlated. The broadcast-show cluster (2 clusters: \emph{heute} with 133 items, \emph{tagesschau} with 508) gives a wider CI ($[+6.56, +10.73]$) because resampling only 2 clusters introduces more variance, but the point estimate ($+9.26$) is close to the query-level. The broadcast-date clustering (297 clusters) matches the query-level CI closely ($[+8.38, +11.45]$). The previous 60-char prefix clustering (627 near-singleton clusters) has been replaced by these interpretable groupings because it provided negligible clustering effect. The donor-pool decomposition analysis of \S\ref{sec:exposure_results} additionally reports separate one-way cluster bootstraps and a joint multiway replicate for the origin contrast, all matching to within $\pm 0.3$ (Appendix~\ref{sec:appendix_decomposition}).

\section{Supplementary Oracle Diagnostics}
\label{sec:appendix_oracle}
\paragraph{Source-conditioned retrieval is legal, oracle-gloss retrieval is leakage.}
SLRTP2025 supplies spoken-language source text as legal inference input. Under pure whole-donor copy, source-text Jaccard scores 23.8 sacreBLEU (0--100 scale, as throughout this appendix), LaBSE 19.2, and multi-SBERT 16.5; the PT-composed source-text condition scores 18.9 (Table~\ref{tab:isomorphic_conditioning}). Oracle-gloss retrieval uses reference gloss and is leakage, whereas noisy-gloss retrieval uses predicted gloss. Thus the strongest legal source-conditioned result is the 23.8 whole-donor stress test, which uses no generator or scaffold.
\paragraph{Disclosure of selection/evaluation reuse.}
The white-box oracle analyses in this appendix reuse the same target reference and evaluator for both donor selection and decoded evaluation. This is an intentional adversarial design: diagonal entries quantify susceptibility to a selector that knows its evaluator, not an independent validation of retrieval quality.

\paragraph{Cross-checkpoint transfer matrix.}
For each query, a selector maximizes its evaluator's teacher-forced likelihood over candidate donor poses; the selected fixed pose is then decoded by every evaluator (token-normalized NLL, sorted donor IDs, \texttt{numpy.argmax} first tie-break). All seven diagonal margins are positive (mean $+0.0150$, 95\% paired-bootstrap CI $[+0.0105,+0.0198]$). Original-selected donors transfer to five of six reconstructions (mean local-GT margin $+0.0074$, $[+0.0016,+0.0135]$); the reverse direction fails in all six cells (reconstruction-selected donors evaluated by Original, mean margin $-0.0455$ $[-0.0554,-0.0358]$). The full 7$\times$7 matrix, unrounded cells, margins, and CIs are in the artifact.



\paragraph{Candidate-pool size.}
The Original-selector stress test shows a strong search-budget effect: at 8, 32, 128, 512 and 2,048 candidates, decoded sacreBLEU remains below local GT by $-12.1$, $-9.2$, $-5.7$, $-3.6$ and $-1.3$ BLEU points respectively; only the full 7,060-donor pool is detectably positive ($+1.7$, $[+0.4,+3.0]$).

\paragraph{Reference-overlap exclusion.}
Masking 1,036 exact-text query--donor pairs before selection: Original-on-Original decoded BLEU is 13.86 versus local GT 12.78 (absolute difference $+1.08$ BLEU points), but the paired relative-margin 95\% CI crosses zero; the frozen joint likelihood-and-interval success rule is met in 0/7 checkpoints. Masking all 19,694 pairs with source-token Jaccard $\geq 0.5$: the score falls to 8.30 ($-4.47$), again 0/7 meet the rule. The full-pool diagonal advantage does not survive as a statistically robust decoded-BLEU result after exact-text exclusion and disappears after high-overlap exclusion.



\paragraph{Noisy-gloss retrieval is split-dependent.}
A 6-layer text$\to$gloss Transformer (dev BLEU-4 13.17) yields noisy-gloss retrieval sacreBLEU 12.7 on PHX-public (strong-NMT), versus 1.6 for a weak NMT (BLEU-4 1.15) and 6.1 for XLM-R embedding retrieval; donor exclusion reduces the strong-NMT value to 10.9 at $\tau{=}0.5$, and the same control reaches 0.4 on the template-disjoint holdout (Table~\ref{tab:holdout_controls}).




\paragraph{PHX-public pure-donor donor-exclusion $\tau$-sweep.}
\label{sec:appendix_public_donor_exclusion}
To test whether the headline reversal survives donor exclusion on the same split, we applied a source-text Jaccard $\tau$-sweep and exact-text-exclusion to the canonical pure text-nearest donor on PHX-public under the original SLRTP2025 BT (Table~\ref{tab:pub_tau_sweep}). The reversal survives exact-text exclusion (23.0 vs GT 12.8, 32/641 donors replaced) and source-Jaccard exclusion up to $\tau{=}0.5$ (17.8 vs GT 12.8); only at $\tau{=}0.3$ and $\tau{=}0.1$ does the retrieval-vs-GT gap turn negative.

\begin{table}[t]
\centering
\small
\caption{PHX-public pure-donor source-text donor-exclusion $\tau$-sweep under the original SLRTP2025 BT (canonical tie-breaking, 12.5\,fps; sacreBLEU on the 0--100 scale). The retrieval-exceeds-GT reversal survives exact-text exclusion and $\tau{=}0.5$ exclusion.}
\label{tab:pub_tau_sweep}
\begin{tabular}{lrrr}
\toprule
Condition & sacreBLEU & $\Delta$ vs GT & Reversal? \\
\midrule
GT baseline                             & 12.8 & ---   & --- \\
Text-nearest pure (no exclusion)        & 23.8 & $+11.0$ & YES \\
Exact-text excluded                     & 23.0 & $+10.2$ & YES \\
$\tau{=}0.7$ (Jaccard $>$ 0.7 excluded) & 21.2 & $+8.4$ & YES \\
$\tau{=}0.5$ (Jaccard $>$ 0.5 excluded) & 17.8 & $+5.0$ & YES \\
$\tau{=}0.3$ (Jaccard $>$ 0.3 excluded) & 8.0 & $-4.8$ & NO \\
$\tau{=}0.1$ (Jaccard $>$ 0.1 excluded) & 0.2 & $-12.5$ & NO \\
Random donor (mean of 20 seeds)          & 0.9 & $-11.6$ & NO \\
\bottomrule
\end{tabular}
\end{table}

As a same-paradigm proxy for the SLRTP2025 winner (whose checkpoints are not public), the noisy-gloss retrieval+LM system shows the same pattern: it survives exclusion up to $\tau{=}0.5$ (10.9, still $6.8\times$ the PT-only baseline and $11.7\times$ the random-donor mean; mean selected-donor Jaccard 0.46) and collapses only at $\tau{=}0.1$.


\section{Template-Density, Holdout Boundary, and Donor-Exclusion Analyses}
\label{sec:appendix_template_boundary}
\label{sec:donor_decomposition}
\label{sec:frozen-template-holdout}
\label{sec:aggressive-donor-audit}

\paragraph{Template density does not gate the reversal.}
Stratifying the 641 test sequences by template density (gloss bigram-Jaccard with nearest train donor), retrieval exceeds GT at all levels: low $+10.7$, mid $+10.1$, high $+12.4$ BLEU points (Table~\ref{tab:template_density_strata}). Density moderates magnitude but does not gate presence.

\begin{table}[h]
\centering
\small
\caption{Template-density stratification (sacreBLEU 0--100). Retrieval exceeds GT at all density levels.}
\label{tab:template_density_strata}
\begin{tabular}{lrrrr}
\toprule
Density stratum & Mean density & GT & Retrieval & Gap \\
\midrule
Low  (n=249)  & 0.152 & 5.9 & 16.6 & $+10.7$ \\
Mid  (n=179)  & 0.254 & 9.5 & 19.7 & $+10.1$ \\
High (n=213)  & 0.510 & 26.1 & 38.6 & $+12.4$ \\
\bottomrule
\end{tabular}
\end{table}

\paragraph{Holdout boundary controls.}
Truncating to 64/32 frames shrinks the gap from $+11.0$ to $+10.5$/$+5.6$ but the reversal persists. Four holdout controls (Table~\ref{tab:holdout_series}) all fail to reproduce the reversal: the same-evaluator seen-holdout gives GT 77.0 (readout), BT-retrained on unseen holdout gives gap $-0.9$, retrained BT on PHX-public gives gap $-0.4$, and cross-fit A/B gives gap $-0.1$.

\begin{table}[h]
\centering
\scriptsize
\caption{Holdout boundary controls (sacreBLEU, 0--100). None reproduces the $+11.0$ reversal.}
\label{tab:holdout_series}
\begin{tabular}{llrrr}
\toprule
Control & Setup & GT & Retrieval & Gap \\
\midrule
Same evaluator & test split (template-overlap) & 12.8 & 23.8 & $+11.0$ \\
Same evaluator & template-holdout (seen in training) & 77.0 & 20.1 & $-56.9$ \\
BT-retrained (unseen holdout) & train-core 6{,}535 & 6.2 & 5.4 & $-0.9$ \\
Retrained BT on PHX-public & unseen to both & 8.7 & 8.3 & $-0.4$ \\
Cross-fit A/B (donor-unseen) & pooled, $n{=}659$ & 4.0 & 3.9 & $-0.1$ \\
\bottomrule
\end{tabular}
\end{table}

\paragraph{Frozen template-holdout with retrained components.}
A deterministic template-holdout (14{,}988 train-core windows, 1{,}538 holdout windows, disjoint by window, source sequence, and exact gloss template) with all components retrained on train-core shows text-nearest retrieval above random ($+2.3$) but far below GT ($-10.8$): the shortcut is real over random but does not exceed GT on the clean split. Donor-exclusion sweeps show the advantage persisting at $\tau{=}0.5$ and collapsing only at $\tau{=}0.1$ (Table~\ref{tab:holdout_controls}).

\begin{table}[h]
\centering
\scriptsize
\caption{Frozen template-holdout controls (1,538 windows, $\alpha{=}0.5$; sacreBLEU 0--100).}
\label{tab:holdout_controls}
\begin{tabular}{lrrr}
\toprule
System & sacreBLEU & Empty \% & Mean Jac. \\
\midrule
GT baseline & 13.10 & 9.95 & --- \\
Random train donor & 0.04 & 11.83 & --- \\
Text-nearest train donor & 2.30 & 10.66 & --- \\
LaBSE / multi-SBERT (donor copy) & 1.73 / 1.51 & 9.2 / 11.3 & --- \\
Noisy-gloss retrieval (strong NMT) & 0.40 & 10.90 & --- \\
Retrieval composition ($\alpha{=}0.5$) & 2.20 & 10.79 & --- \\
\midrule
$\tau{=}0.1$ / 0.3 / 0.5 / 0.7 exclusion & 0.04 / 1.05 / 1.93 / 2.01 & 10.0--10.5 & 0.10/0.29/0.40/0.42 \\
BM25 / TF-IDF nearest & 1.75 / 2.36 & 10.8 / 8.8 & --- \\
\bottomrule
\end{tabular}
\end{table}

The SLRTP2025 test vocabulary is nearly closed (95.4\% gloss types in train); nearest donors share mean gloss unigram containment 0.557 and bigram-Jaccard 0.302, but Progressive Transformer outputs show no donor-proximity effect ($\rho{=}{-}0.06$, $p{=}0.12$), versus $+0.29$ for retrieval.

\section{DGS Human Evaluation Details}
\label{sec:appendix_human_eval}
This appendix gives the full data for the blind DGS human evaluation feasibility study described in \S\ref{sec:human_eval_feasibility}. The 100 video files and the UUID-to-system manifest were lost after the evaluation was completed, so per-system analyses cannot be made directly. We instead report the rank-based profile and a 120-permutation exploration that bounds what could be inferred if the mapping were recovered.

\paragraph{Per-rank profile.}
For each of the 20 target sentences, the 5 videos were ranked by mean semantic-adequacy across the 30 raters (rank 1 = highest, rank 5 = lowest). Per-rank aggregate (mean across the 20 targets):

\begin{table}[h]
\centering
\scriptsize
\caption{Per-rank profile across 20 target sentences (30 raters per video; 5 videos per target). Ranks are defined by mean semantic-adequacy within each target; the table is descriptive only.}
\label{tab:appendix_human_eval_rank}
\begin{tabular}{lrrrrr}
\toprule
Rank & $n$ & Semantic mean (SD) & Intell.\ mean & Naturalness mean & Semantic 95\% CI \\
\midrule
1 & 20 & 4.48 (0.16) & 3.82 & 3.90 & $[4.41, 4.55]$ \\
2 & 20 & 3.62 (0.67) & 3.75 & 3.48 & $[3.31, 3.94]$ \\
3 & 20 & 2.89 (0.88) & 3.38 & 3.34 & $[2.48, 3.30]$ \\
4 & 20 & 2.28 (0.67) & 3.38 & 3.54 & $[1.97, 2.59]$ \\
5 & 20 & 1.83 (0.43) & 3.59 & 3.71 & $[1.63, 2.03]$ \\
\bottomrule
\end{tabular}
\end{table}

\paragraph{Inter-rater agreement.}
Krippendorff's $\alpha$ (ordinal) on the complete rater $\times$ video matrix ($30 \times 100$):
\begin{itemize}
\item Semantic adequacy: $\alpha = 0.543$
\item Intelligibility: $\alpha = 0.480$
\item Naturalness: $\alpha = 0.490$
\end{itemize}
These values are in the moderate-agreement range (Krippendorff's conventional threshold for ``substantial'' agreement is $\alpha \geq 0.667$; moderate is $0.44 \leq \alpha < 0.667$). The semantic-adequacy scale, which is the most directly relevant for the BT-reversal question, has the highest agreement of the three.

\paragraph{Permutation exploration (no ground-truth mapping).}
The 100 videos were rendered from 5 canonical systems (GT, PT-v1, TN-PURE-v1, TN-PTCOMP-v1, random donor) on 20 target sentences, with system labels replaced by UUIDs for blind rating. Without the lost UUID-to-system manifest, no per-system analysis is possible. For transparency, we explored all 120 permutations of the 5 systems against the 5 DGS-rank positions. The Spearman correlation between corpus BT-BLEU (paper-reported: GT 12.78, PT 1.59, PURE 23.79, COMP 18.86, random 0.90) and per-rank mean DGS-semantic-adequacy ranges from $\rho = -0.97$ (most inverted) to $\rho = +0.98$ (most concordant); without the manifest we cannot identify which permutation is correct. We therefore make no per-system claim from the human evaluation.

\paragraph{What is independent of the mapping.}
The following findings hold for any UUID-to-system permutation, because they are rank-based:
\begin{enumerate}
\item DGS raters agree moderately on semantic adequacy ($\alpha = 0.54$), confirming that semantic judgment is possible without the BT-decoded text.
\item The 5 videos per target span a wide DGS-semantic range (top rank 4.48 vs bottom rank 1.83, $\Delta = 2.65$ on a 1--5 scale).
\item Because ranks are defined by semantic adequacy, the rank profile is descriptive rather than an independent test. Naturalness and intelligibility show less variation across these semantic rank bins, but the lost manifest prevents any system-level attribution.
\end{enumerate}
The mapping-independent result is moderate inter-rater agreement on semantic adequacy ($\alpha = 0.54$). The rank profile (top vs.\ bottom video) is descriptive and cannot support system attribution or a correlation analysis against decoded BT-text metrics: top and bottom ranks differ by construction because they are defined by the semantic-adequacy scores, and after the manifest was lost we cannot test whether the rank order is concordant or inverted with corpus BT-BLEU.

\paragraph{Limitations of this human evaluation.}
First, the lost manifest prevents per-system claims; a re-rendering of the 100 videos from the canonical pose tensors and a re-rating would restore the per-system analysis. Second, DGS-fluent raters were not asked to judge spatial grammar, indexation, or non-manual markers (mouth shapes, head tilt, eye gaze); the semantic-adequacy scale is a holistic target-meaning match. Third, we did not collect rater demographics (age, L1/L2 status, signing experience) separately; recruitment was through DGS-fluent channels. Fourth, the target-sentence set is 20 German weather sentences from PHX-public; results may not generalize beyond this domain.

\section{Transcript-Slot Diagnostic Details}
\label{sec:appendix_slot}
This appendix documents the rule-based German weather slot extractor referenced in \S\ref{sec:slot_f1}. The extractor identifies four slot families (number/temperature, time, place, weather-event) using a frozen lexicon and normalization rules. On a 50-reference audit (seed 20260729), per-family precision was 0.91--1.00 and recall 0.75--1.00 (overall precision 0.97, recall 0.86); errors are dominated by recall misses from inflections, compounds, and lexicon gaps. A stratified audit of 50 decoded hypotheses (25 GT-decode, 25 PURE-decode, seed 20260730) showed precision nearly equal (0.96 vs 0.97) but recall lower on PURE-decode (0.83 vs 0.91), because donor-transcript-style outputs use more inflected verbs and compounds (\emph{regnet}, \emph{unwetterwarnungen}) that the lexicon misses. The measured PURE slot-F1 is thus attenuated more than GT's, so the PURE${>}$GT contrast is, if anything, underestimated. The audit used a single auditor; inter-annotator agreement is unavailable. The full lexicon, normalization rules, and per-family precision/recall are in the artifact.

\section{Competence Extrapolation Details (Exploratory)}
\label{sec:appendix_extrapolation}
This appendix gives the full OLS/GP extrapolation analysis referenced in \S\ref{sec:extrapolation}. We emphasize that this is exploratory and the prediction intervals are descriptive visual guides, not calibrated inferential tests, because the 43 checkpoints share data, code, and hyperparameters (non-IID) and the target point (dev BLEU-4 $=13.38$) lies outside the observed range.

\paragraph{Method.}
We fit the gap--competence dose--response over all 43 evaluated checkpoints spanning dev BLEU-4 $[4.5, 10.9]$ using three models: OLS linear, OLS quadratic, and a Gaussian Process with RBF + white kernel. Each model predicts the PURE--GT gap at dev BLEU-4 $=13.38$.

\paragraph{Results.}
Table~\ref{tab:extrapolation} reports the predictions and 95\% PIs at dev BLEU-4 $=13.38$. All predictions are negative, far below the observed $+11.01$. Table~\ref{tab:loo_extrapolation} reports a leave-one-family-out cross-validation; the OLS quadratic prediction is the most sensitive (flips from $-13.5$ to $+1.3$ when the 15 distillation students are removed), while the GP upper bound stays below $+2.0$ across all five folds.

\begin{table}[h]
\centering
\small
\caption{Competence extrapolation: gap predictions at dev BLEU-4 $=13.38$ from three fits over 43 checkpoints (dev BLEU-4 4.5--10.9). All predictions are far below the observed $+11.01$. PIs are \emph{descriptive} (non-IID checkpoints); they should not be read as calibrated uncertainty intervals.}
\label{tab:extrapolation}
\begin{tabular}{lrrr}
\toprule
Fit & R$^2$ / LML & Prediction at 13.38 & 95\% PI (descriptive) \\
\midrule
OLS linear & 0.58 & $-6.4$ & $[-9.5, -3.3]$ \\
OLS quadratic & 0.78 & $-13.5$ & $[-16.9, -10.1]$ \\
Gaussian Process (RBF + white) & $-29.5$ & $-2.1$ & $[-6.3, +2.0]$ \\
\midrule
Released evaluator (observed) & --- & $+11.01$ & $[+9.6, +12.4]$ \\
\bottomrule
\end{tabular}
\end{table}

\begin{table}[h]
\centering
\scriptsize
\caption{Leave-one-family-out cross-validation of the competence extrapolation. Each row refits all three models after removing one checkpoint family. The GP upper bound stays below $+2.0$ across all folds; OLS quadratic is the most sensitive (flips sign when distillation is removed). All intervals descriptive.}
\label{tab:loo_extrapolation}
\begin{tabular}{lrrrrrr}
\toprule
Family removed & $n$ remaining & OLS-lin pred & OLS-quad pred & GP pred & GP upper & OLS-lin R$^2$ \\
\midrule
None (full) & 41 & $-6.4$ & $-13.5$ & $-2.1$ & $+2.0$ & 0.58 \\
Reconstruction (14) & 27 & $-6.6$ & $-15.4$ & $-2.9$ & $+2.0$ & 0.57 \\
Distillation (13) & 28 & $-1.1$ & $+1.3$ & $-0.4$ & $+0.5$ & 0.14 \\
Ladder (4) & 37 & $-8.1$ & $-18.0$ & $-2.3$ & $+1.8$ & 0.66 \\
Config-faithful (4) & 37 & $-6.7$ & $-13.2$ & $-2.4$ & $+1.9$ & 0.63 \\
Rescue (2) & 39 & $-6.8$ & $-13.7$ & $-2.1$ & $+1.7$ & 0.65 \\
\bottomrule
\end{tabular}
\end{table}

\paragraph{Caveats.}
First, the GP upper bound ($+2.0$) partly reflects RBF kernel prior shrinkage toward the mean in the extrapolation region (sparse data beyond dev BLEU-4 10.9), not purely data evidence. Second, the OLS quadratic fit is sensitive to the distillation family. Third, the 43 checkpoints are not IID. We therefore report these fits only as descriptive visualizations and explicitly avoid claiming that ``competence alone is unlikely to explain the reversal'' on the basis of these fits. The unobserved 10.9--13.4 interval could in principle harbor non-linear behavior that the present experiments cannot rule out.

\section{Supplementary CSL-Daily Boundary Control}
\label{sec:csl-daily}
CSL-Daily (Chinese Sign Language) serves as a non-template-domain control, but its recognizer is too weak for robust claims. Under Mode B (1{,}176 sequences), text-nearest donor copy reaches only 1.6 sacreBLEU vs GT 26.7 (Table~\ref{tab:csl_donor_exclusion}). A readout-ratio check (64 train + 64 test windows) gives train$\approx$test BLEU-1 (0.0019 vs 0.0021, ratio $\sim$0.9), far below PHX's 6.1: the readout signature does not generalize beyond the released PHX evaluator.

\begin{table}[h]
\centering
\small
\caption{CSL-Daily Mode B donor-exclusion (1{,}176 sequences; 0--100 scale).}
\label{tab:csl_donor_exclusion}
\begin{tabular}{lrr}
\toprule
Condition & sacreBLEU & BLEU-1 \\
\midrule
GT baseline & 26.7 & 79.1 \\
Text-nearest donor & 1.6 & 43.5 \\
Random donor & 0.0 & 3.9 \\
$\tau{=}0.5$ exclusion & 1.0 & 42.6 \\
BM25 retrieval & 1.5 & 39.2 \\
TF-IDF retrieval & 1.2 & 38.5 \\
\bottomrule
\end{tabular}
\end{table}

\section{Leakage Sanity Checks for Train-Pool Readout}
\label{sec:appendix_leakage}
This appendix documents the five sanity checks referenced in \S\ref{sec:leakage_checks} that rule out implementation-level artifacts as explanations for the released evaluator's 78.8 BLEU / 70.7\% EM train-pool readout.

\paragraph{Experiment A: Code-path audit.}
The \texttt{back\_translate()} function in \texttt{src/models/back\_translate.py} (lines 362--426) calls \texttt{model.run\_batch(translation\_beam\_size=3, translation\_beam\_alpha=-1)}, which invokes \texttt{bt\_search.beam\_search()}. The decoder input is seeded with \texttt{<bos>} only and extended autoregressively with the model's own previously generated tokens. The reference text is used exclusively in the post-hoc BLEU/EM scoring step; it never enters the model forward pass. We verified this by code inspection and by confirming that the same model loaded in a fresh Python process produces identical outputs.

\paragraph{Experiment B: Pose--text permutation.}
We decoded the first 1{,}000 training items with beam=3 (BLEU 78.61, EM 70.1\%), then permuted the reference labels across items to break pose--text correspondence. Three random permutations yield:
\begin{table}[h]
\centering
\small
\caption{Label permutation test (1{,}000 items, reference labels shuffled \emph{after} decoding). BLEU collapses from 78.6 to $\sim$1.0, confirming output--reference alignment. This test alone cannot distinguish pose-conditioned output from ID/cache-based retrieval; the input-side pose permutation below addresses that.}
\label{tab:permutation_test}
\begin{tabular}{lrr}
\toprule
Condition & BLEU & EM \\
\midrule
Original (correct pairing) & 78.61 & 0.701 \\
Permutation 0 & 1.06 & 0.000 \\
Permutation 1 & 1.30 & 0.000 \\
Permutation 2 & 0.90 & 0.002 \\
\bottomrule
\end{tabular}
\end{table}

\paragraph{Experiment B$'$: Input-side pose permutation with output-equivariance.}
\label{sec:input_permutation}
The label permutation above shuffles reference labels \emph{after} decoding, so it cannot exclude ID-based or cache-based retrieval (a cache lookup by ID would also produce low BLEU when labels are shuffled). To properly test pose conditioning, we held item IDs and text metadata fixed, shuffled the \emph{pose tensors} across 100 training items \emph{before decoding}, and performed output-equivariance verification.
\begin{table}[h]
\centering
\small
\caption{Input-side pose permutation with output-equivariance (100 items). Pose tensors are shuffled \emph{before} decoding while IDs stay fixed. The equivariance column verifies $E(\text{pose}_j, \text{slot}_i) = E(\text{pose}_j, \text{slot}_j)$.}
\label{tab:input_permutation}
\begin{tabular}{lrr}
\toprule
Condition & BLEU vs slot ref & Detail \\
\midrule
Original decode vs own reference & 75.62 & baseline \\
Shuffled poses vs slot reference & 0.95 & BLEU collapses \\
Shuffled poses vs donor reference & 75.62 & identical to original \\
\midrule
Output-equivariance match rate & --- & 100/100 (100\%) \\
Donor-length match rate & --- & 100/100 (100\%) \\
Slot-length match rate & --- & 5/100 (5\%) \\
\bottomrule
\end{tabular}
\end{table}
For all 100 items, the hypothesis produced when pose$_j$ is decoded in slot $i$'s position exactly matches the hypothesis originally produced for pose$_j$ in its own slot $j$ (100\% exact string match). Output length follows the pose tensor, not the slot ID (100\% donor-length match vs.\ 5\% slot-length match). This rules out a pose-ignorant pure ID/cache lookup under the tested decoding path and confirms sensitivity to the substituted pose tensor. Additional degenerate-pose controls (zero, Gaussian noise, single repeated pose for 200 items) all collapse BLEU to $<$1, and reversed batch order reproduces the original BLEU exactly, confirming per-item determinism.

\paragraph{Experiment C: Determinism.}
Beam search with no stochastic sampling produces deterministic outputs. We record SHA-256 hashes of all per-item decoded hypotheses in the artifact for cross-process verification.

\paragraph{Experiment D: Stratified EM.}
Exact-match rate is high and broadly distributed, not concentrated in a single subgroup:
\begin{table}[h]
\centering
\scriptsize
\caption{Stratified exact-match rate on the full 7{,}060-item training pool under beam=3 free-decode.}
\label{tab:stratified_em}
\begin{tabular}{lrr}
\toprule
Stratum & $n$ & EM \\
\midrule
\textbf{By signer} & & \\
Signer01 & 1{,}852 & 0.683 \\
Signer02 & 83 & 0.590 \\
Signer03 & 579 & 0.736 \\
Signer04 & 1{,}049 & 0.748 \\
Signer05 & 1{,}623 & 0.701 \\
Signer06 & 35 & 0.714 \\
Signer07 & 762 & 0.696 \\
Signer08 & 844 & 0.715 \\
Signer09 & 233 & 0.734 \\
\midrule
\textbf{By reference length} & & \\
Short ($\leq$5 words) & 203 & 0.837 \\
Medium (6--10) & 1{,}319 & 0.685 \\
Long (11--20) & 4{,}335 & 0.714 \\
Very long ($>$20) & 1{,}203 & 0.686 \\
\midrule
\textbf{By text repetition in train} & & \\
Unique (1 occurrence) & 6{,}777 & 0.702 \\
Low (2--3) & 57 & 0.737 \\
Mid (4--10) & 76 & 0.684 \\
High ($>$10) & 150 & 0.940 \\
\bottomrule
\end{tabular}
\end{table}
Even unique training texts (appearing exactly once) achieve 70.2\% EM, confirming the readout is not driven by template repetition alone. Of the 7{,}060 hypotheses, 74.8\% match some training text exactly (70\% match exactly one---their correct paired reference).

\paragraph{Experiment E: Same-signer held-out control.}
We decoded 500 test-split poses (unseen during training) from the same 9 signers present in the training pool. This is a same-signer held-out control, not a fully matched membership control: the test items are not matched on length, text-similarity, broadcast date, or template density, so the train--test gap reflects a combination of training-pool exposure and distributional shift.
\begin{table}[h]
\centering
\small
\caption{Same-signer held-out control: train (seen) vs.\ test (unseen) free-decode. The 67-point BLEU gap is consistent with training-pool exposure contributing to the readout, though the lack of full covariate matching means we cannot quantify how much of the gap is exposure vs.\ distributional shift.}
\label{tab:membership_control}
\begin{tabular}{lrrrr}
\toprule
Split & $n$ & BLEU & EM & Shared signers \\
\midrule
Train (seen) & 7{,}060 & 78.82 & 0.707 & 9 \\
Test (unseen) & 500 & 12.19 & 0.032 & 9 (same) \\
\bottomrule
\end{tabular}
\end{table}
The same signers appear in both splits. The 67-point BLEU gap between seen and unseen items is consistent with training-pool exposure contributing to the readout, but without full covariate matching (length, text-similarity, template density), we cannot attribute the gap solely to exposure vs.\ distributional shift.


\bibliography{references}

\end{document}
